\chapter{1990:Elias James Corey}

\label{chap:chemistry1990}

The Nobel Prize in Chemistry for the year 1990 was awarded to Elias James Corey.

\textbf{Motivation:}

for his development of the theory and methodology of organic synthesis

\section{Elias James Corey}

\subsection*{Early Life and Education}

Elias James Corey was born in 1928 in Methuen, Massachusetts, USA.

\subsection*{Career and Major Research}

Corey studied at the Massachusetts Institute of Technology (B.S., 1948; Ph.D., 1951). He was a professor at the University of Illinois from 1951 to 1959, and then joined Harvard University, where he spent his career.

\subsection*{Nobel Prize-winning Work}

The work for which Elias James Corey was awarded the Nobel Prize in Chemistry in 1990 was described by the Nobel Committee as: "for his development of the theory and methodology of organic synthesis".

He developed retrosynthetic analysis, a systematic method for planning complex organic syntheses by disconnecting target molecules into simpler precursor structures.

\subsection*{Significance and Impact}

Retrosynthetic analysis became the standard conceptual framework of organic synthesis and guided the synthesis of hundreds of natural products.

\subsection*{Later Career and Honors}

Corey continued teaching and research at Harvard University, where he is professor emeritus. He received the 1990 Nobel Prize in Chemistry.

\subsection*{Legacy}

His methodology, which includes the Corey-Bakshi-Shibata asymmetric reduction, remains central to synthetic organic chemistry.

\subsection*{Selected Publications}

"The Logic of Chemical Synthesis" (Wiley, 1989)

"Computer-Assisted Design of Complex Organic Syntheses" (Science, 1969)

\clearpage

\section*{Chapter Summary}

The 1990 Nobel Prize in Chemistry recognized the theory and methodology of organic synthesis. Corey's retrosynthetic analysis gave chemists a rational and teachable strategy for constructing complex molecules.
